\section*{Bloque C. Casos con implicancias regulatorias y sector saneamiento}

\noindent\textit{Preguntas 66 a 90. Evalúan la aplicación al marco de SUNASS y saneamiento peruano: Valores Máximos Admisibles (VMA), mecanismos MERESE en cuencas, agua como bien meritorio, subsidios focalizados (SISFOH), Agua No Facturada (ANF), JASS y sobreexplotación de acuíferos.}

\q{Los Mecanismos de Retribución por Servicios Ecosistémicos Hídricos (MERESE), incorporados por SUNASS en las tarifas de agua, representan la aplicación práctica de:}
\begin{opts}
  \item un subsidio ciego a las planillas administrativas de las EPS.
  \item una solución coaseana e instrumento económico donde los usuarios urbanos de agua retribuyen a las comunidades de la cuenca alta para financiar la conservación de bofedales y la regulación hídrica.
  \item un impuesto arancelario a la importación de tecnología hídrica.
  \item una privatización forzosa de los ríos y lagunas del país.
\end{opts}

\q{Los Valores Máximos Admisibles (VMA) regulados por el Ministerio de Vivienda y fiscalizados por las EPS y SUNASS tienen como objetivo económico primordial:}
\begin{opts}
  \item fijar una tasa de interés preferencial para préstamos a microempresas.
  \item prohibir la actividad industrial y comercial en todas las ciudades del Perú.
  \item controlar la externalidad negativa que causan las descargas de aguas residuales no domésticas con alta carga contaminante en las redes de alcantarillado, evitando el deterioro prematuro de tuberías y PTAR.
  \item eliminar el cobro de alcantarillado a los usuarios residenciales.
\end{opts}

\q{El pago adicional por exceso de concentración de VMA que una EPS factura a una industria contaminante opera económicamente como:}
\begin{opts}
  \item una donación caritativa deducible de impuestos.
  \item un monopolio bilateral de comercialización de aguas negras.
  \item una tasa de rentabilidad libre de riesgo según el CAPM.
  \item un instrumento económico pigouviano que penaliza el daño marginal y genera incentivos para que la industria instale su propio sistema de pretratamiento.
\end{opts}

\q{En la elaboración del Plan Maestro Optimizado (PMO) de una EPS para fijar las tarifas quinquenales ante SUNASS, existe una asimetría de información de tipo selección adversa porque:}
\begin{opts}
  \item la EPS tiene información privada sobre su verdadera estructura de costos de inversión y operación, teniendo incentivos estratégicos a sobreestimar costos para obtener tarifas más altas.
  \item SUNASS conoce mejor la geología de la cuenca que los ingenieros locales de la EPS.
  \item los usuarios residenciales ocultan su nivel de consumo mensual a la EPS.
  \item el Banco Mundial prohíbe la auditoría técnica de los expedientes tarifarios.
\end{opts}

\q{En el marco de la regulación económica en el Perú, el agua potable y saneamiento se considera un 'Bien Meritorio' debido a que:}
\begin{opts}
  \item no requiere inversión en infraestructura física para su potabilización.
  \item el consumo de agua segura genera externalidades positivas de salud pública (prevención de anemia y desnutrición) que los hogares de menores ingresos podrían desatender si no existiera una tarifa subsidiada.
  \item es un bien club privado administrado exclusivamente por corporaciones transnacionales.
  \item su demanda es perfectamente elástica respecto al ingreso monetario.
\end{opts}

\q{El mecanismo de subsidios cruzados focalizados mediante el Sistema de Focalización de Hogares (SISFOH) aplicado por SUNASS busca:}
\begin{opts}
  \item eliminar el uso de micromedidores en toda la provincia.
  \item otorgar agua gratuita a todas las industrias y comercios de la ciudad.
  \item dirigir el subsidio tarifario únicamente a los hogares en condición de pobreza y pobreza extrema, evitando las fugas del subsidio hacia estratos socioeconómicos no vulnerables.
  \item financiar campañas publicitarias de la empresa prestadora.
\end{opts}

\q{En las Juntas Administradoras de Servicios de Saneamiento (JASS) en el ámbito rural peruano, la alta morosidad en el pago de la 'cuota familiar' suele ser un caso emblemático de:}
\begin{opts}
  \item discriminación de precios de tercer grado.
  \item competencia ruinosa a la Bertrand.
  \item colusión tácita entre constructores de reservorios.
  \item comportamiento de polizón (free-rider), donde los comuneros aprovechan la dificultad social de cortar el servicio para eludir el pago del mantenimiento preventivo del sistema.
\end{opts}

\q{Los elevados niveles de Agua No Facturada (ANF), que en muchas EPS peruanas superan el 40\%, representan desde la perspectiva del regulador un problema de:}
\begin{opts}
  \item riesgo moral operativo y deficiencia de gestión, donde la falta de incentivos de alto poder permite desidia en el mantenimiento de redes y tolerancia al clandestinaje.
  \item selección adversa en la importación de macro y micromedidores.
  \item externalidad de red positiva que favorece a los asentamientos informales.
  \item monopolio bilateral con la Autoridad Nacional del Agua.
\end{opts}

\q{La crítica sobreexplotación del acuífero subterráneo del valle de Ica debido a la masiva proliferación de pozos tubulares para agroexportación es una manifestación directa de:}
\begin{opts}
  \item un equilibrio de Lindahl perfectamente balanceado.
  \item la Tragedia de los Comunes, donde la falta de medición estricta de caudales y fiscalización de derechos induce a cada usuario a extraer agua antes de que se agote la napa freática.
  \item un cártel colusorio coordinado por la Cámara de Comercio de Ica.
  \item la aplicación exitosa de la condición de provisión de Samuelson.
\end{opts}

\q{En los procesos de contratación de obras de saneamiento bajo el sistema Invierte.pe, cuando un contratista presenta una oferta económica artificialmente baja confiando en que posteriormente recuperará rentabilidad mediante 'adendas y ampliaciones de plazo', se produce:}
\begin{opts}
  \item un equilibrio agrupador con bienestar social máximo.
  \item un mecanismo de revelación de preferencias eficiente a la Groves-Ledyard.
  \item un fenómeno de selección adversa y oportunismo postcontractual facilitado por deficiencias del expediente técnico.
  \item una solución coaseana de bajo costo de transacción.
\end{opts}

\q{En un contrato de concesión para una Planta de Tratamiento de Aguas Residuales (PTAR) bajo modalidad APP, si el Estado remunera al operador exclusivamente mediante un pago fijo por disponibilidad (CAP) sin fiscalizar la calidad del agua de salida (LMP):}
\begin{opts}
  \item el operador tratará el agua con estándares superiores a los exigidos.
  \item se garantiza la sostenibilidad financiera de la cuenca hídrica.
  \item el precio de los permisos de emisión caerá a cero.
  \item se genera un severo riesgo moral, donde el concesionario reduce costos operativos apagando los sistemas de aireación o escatimando reactivos químicos, vertiendo agua mal tratada al cuerpo receptor.
\end{opts}

\q{El mecanismo de 'Benchmarking Regulatorio' aplicado por SUNASS para comparar la eficiencia técnica y de costos entre las 50 EPS del país utiliza modelos econométricos para:}
\begin{opts}
  \item estimar una frontera de costos eficientes (DEA o SFA) que permita recortar las rentas informacionales de las empresas y fijar metas de costos exigentes en las tarifas.
  \item reemplazar las audiencias públicas de rendición de cuentas.
  \item otorgar premios monetarios directos a los alcaldes provinciales.
  \item unificar el precio del metro cúbico de agua en todo el territorio nacional.
\end{opts}

\q{La provisión de agua potable con desinfección mediante cloro residual genera una externalidad positiva masiva sobre el sistema público de salud porque:}
\begin{opts}
  \item permite a las EPS privatizar sus activos operativos.
  \item previene brotes de cólera, tifoidea y diarreas agudas, reduciendo la saturación hospitalaria y los días de trabajo y escuela perdidos por enfermedad.
  \item incrementa la inflación subyacente reportada por el BCRP.
  \item reduce inmediatamente el precio de los alimentos procesados en supermercados.
\end{opts}

\q{La instalación de micromedidores en conexiones domésticas que anteriormente pagaban tarifa asignada (consumo no medido) reduce el consumo promedio por hogar entre un 20\% y 30\%. Este impacto económico se debe a que:}
\begin{opts}
  \item la EPS raciona el agua mediante cortes de suministro obligatorios.
  \item se incrementan los costos fijos de la factura mensual.
  \item el medidor hace que el costo marginal percibido por cada litro consumido pase de cero a la tarifa volumétrica regulada, eliminando el riesgo moral del desperdicio.
  \item el medidor altera los gustos y preferencias de los miembros del hogar.
\end{opts}

\q{El 'Benchmarking Regulatorio y Ranking de Desempeño' publicado anualmente por SUNASS actúa en los mercados de servicios públicos como un instrumento de:}
\begin{opts}
  \item discriminación perfecta de precios de primer grado.
  \item subsidio ciego a las EPS ubicadas en los últimos lugares del ranking.
  \item regulación basada en costos históricos contables.
  \item señalización y reducción de asimetría de información que permite a los usuarios, autoridades y entes rectores evaluar la calidad de gestión de los directores y gerentes de las EPS.
\end{opts}

\q{Los conflictos socioambientales por el agua en zonas de cabecera de cuenca entre comunidades agrícolas y proyectos mineros en el Perú reflejan frecuentemente el fracaso de una solución coaseana debido a:}
\begin{opts}
  \item la falta de definición unívoca de los derechos de propiedad consuetudinarios, asimetrías de poder negociador y elevados costos de desconfianza mutua.
  \item que el agua es un bien privado perfectamente transferible en bolsa.
  \item la existencia de costos de transacción insignificantes entre las partes.
  \item la intervención obligatoria del FMI en todos los acuerdos comunales.
\end{opts}

\q{En el marco de la Ley de Mecanismos de Retribución por Servicios Ecosistémicos (Ley 30215), el acuerdo entre una EPS urbana y una comunidad campesina altoandina para restaurar amunas y zanjas de infiltración opera como:}
\begin{opts}
  \item un impuesto distrital sobre el alcantarillado pluvial.
  \item un contrato bilateral privado coaseano canalizado y legitimado a través del estudio tarifario aprobado por SUNASS.
  \item un monopolio de compra de terrenos forestales.
  \item un arbitrio municipal no regulado.
\end{opts}

\q{En condominios o quintas donde existe un único medidor de agua general y la factura se divide equitativamente entre los departamentos sin medidores individuales, surge:}
\begin{opts}
  \item la eliminación total de las pérdidas comerciales de agua.
  \item un consumo menor al óptimo debido a la competencia perfecta.
  \item un problema de polizón (free-rider) y riesgo moral, pues cada familia asume solo una fracción $1/N$ del costo de su consumo adicional, induciendo a un consumo mayor al socialmente óptimo.
  \item una solución eficiente de Lindahl.
\end{opts}

\q{La ampliación de la red pública de alcantarillado a un vecindario genera externalidades de red y de salud para los predios colindantes que aún usan letrinas o pozos sépticos debido a que:}
\begin{opts}
  \item obliga a la municipalidad a clausurar todas las pistas y veredas.
  \item disminuye el valor predial de las viviendas conectadas.
  \item reduce la presión de agua de la tubería matriz.
  \item la erradicación de pozos sépticos filtra menos patógenos a la napa freática compartida y elimina focos de vectores infecciosos en todo el entorno vecinal.
\end{opts}

\q{Los fondos de reserva para Gestión del Riesgo de Desastres (GRD) y Adaptación al Cambio Climático que SUNASS incluye en la tarifa de agua actúan como:}
\begin{opts}
  \item un seguro financiero precautorio internalizado en los costos corrientes para responder ante eventos extremos (como huaicos causados por El Niño) y garantizar la continuidad del servicio.
  \item un préstamo blando no reembolsable al gobierno regional.
  \item un dividendo en efectivo para los accionistas de la EPS.
  \item una donación deducible del impuesto a la renta corporativa.
\end{opts}

\q{El Régimen de Apoyo Transitorio (RAT) del OTASS para reflotar EPS insolventes puede generar un problema de riesgo moral financiero en las municipalidades provinciales si:}
\begin{opts}
  \item el personal técnico de la empresa renuncia en su totalidad.
  \item las autoridades locales anticipan que ante cualquier quiebra o mala gestión, el gobierno nacional asumirá las deudas e inyectará capital financiero sin costo para el municipio.
  \item SUNASS eleva las tarifas de agua en un 50\%.
  \item los alcaldes gestionan la EPS con máxima pulcritud técnica.
\end{opts}

\q{El vertimiento ilegal de aceites vegetales usados por restaurantes de pollos a la brasa en las alcantarilladas de la ciudad es una externalidad negativa que genera:}
\begin{opts}
  \item la reducción instantánea de la tarifa media de agua.
  \item la privatización obligatoria del servicio de alcantarillado.
  \item obstrucción de tuberías por bolas de grasa, atoros y desbordes de aguas residuales en la vía pública, elevando el costo de mantenimiento que pagan todos los usuarios.
  \item un enriquecimiento injustificado de los fabricantes de cloro.
\end{opts}

\q{Un sistema de alcantarillado condominal operado de forma cerrada por un grupo de vecinos en una ladera urbana se aproxima a la categoría microeconómica de:}
\begin{opts}
  \item un recurso de libre acceso sin costo de exclusión.
  \item un monopolio de oferta perfectamente competitivo.
  \item un bien público puro de alcance planetario.
  \item un Bien Club, donde los beneficios son compartidos y no rivales entre los miembros del condominio pero se puede excluir a quienes no aporten a la red interna.
\end{opts}

\q{La regulación tarifaria de SUNASS basada en metas de gestión y penalizaciones económicas sobre las fórmulas tarifarias frente a incumplimientos de calidad busca:}
\begin{opts}
  \item alinear las funciones de utilidad de la gerencia de la EPS con el bienestar social, combatiendo la ineficiencia por riesgo moral regulatorio.
  \item reducir la recaudación fiscal del Ministerio de Economía y Finanzas.
  \item eliminar el riesgo de quiebra de contratistas de obras.
  \item permitir que las EPS operen sin presentar planes de inversión.
\end{opts}

\q{En proyectos de saneamiento rural donde se aplica el Método de Valoración Contingente para estimar la Disposición a Pagar (DAP) por agua intradomiciliaria, el sesgo estratégico de los encuestados consiste en:}
\begin{opts}
  \item revelar siempre con exactitud matemática el ingreso del hogar.
  \item subdeclarar la disposición a pagar si creen que de su respuesta dependerá la tarifa que les cobrarán, o sobredeclararla si creen que de ello depende que el Estado construya la obra gratuita.
  \item negarse a participar en toda consulta electoral.
  \item exigir que el agua sea exportada a países vecinos.
\end{opts}

